\documentclass[11pt]{article}
\usepackage{amsmath,amssymb,amsthm,hyperref}

\title{QPI and Twin Primes: A Geometric Reframing (Conjectures + Evidence)}
\author{Christopher Michael LeBlanc Sr.}
\date{\today}

\newtheorem{conjecture}{Conjecture}
\newtheorem{definition}{Definition}
\newtheorem{remark}{Remark}

\begin{document}
\maketitle

\begin{abstract}
This paper separates proven structural equivalences from open problems.
It frames twin primes via midpoint predicates (Yellow/Red) and records
finite computational evidence, while leaving all infinite claims explicitly conjectural.
\end{abstract}

\section{Midpoint framing}
Define Yellow and Red as in Paper 1. For $p>5$, any twin prime pair $(p,p+2)$
has midpoint $m=p+1$ divisible by $3$ and ending in $0,2,$ or $8$; thus $m$ lies in the Yellow set.
The pair is twin iff $m$ is Red.

\section{What remains open}
\begin{conjecture}[Twin Prime Conjecture]
There exist infinitely many twin primes.
\end{conjecture}

\begin{remark}
Any argument of the form ``Yellow is infinite, therefore Red is infinite'' must introduce
a new mechanism that forces infinitely many Yellow midpoints to have both prime neighbors.
No such mechanism is proven here.
\end{remark}

\section{Empirical evidence}
Include receipts up to the maximum validated bound, including SHA-256 hashes and hardware specs.

\end{document}
